\appendix
\section{Pre-Registration and Adversarial Audit Protocol}
\label{sec:appendix-methodology}

Every honest-negative-result claim in \S\ref{sec:discussion} rests on a
process claim as much as a numerical one: that a ``dead'' cell was actually
tested hard enough to trust, and that a ``positive'' cell is not a
shortcut in disguise. This appendix makes that process auditable rather
than asserted, since it is the paper's main methodological differentiator
from a result reported without it.

\subsection{Pre-registration}
\label{sec:appendix-prereg}

Each of the three experiments (Task D, Task E, Stage 0) has its own
pre-registration document -- \texttt{TASK\_D\_PREREGISTRATION.md},
\texttt{NEXT\_EXPERIMENT\_DESIGN.md}, \texttt{STAGE0\_DESIGN.md} in the
repository (\S\ref{sec:repro}) -- written and committed \emph{before} any
GPU time was spent on that experiment. Each fixes the hypothesis, the
architecture and controls, the primary metric, and CONFIRM/FALSIFY
decision bands in advance. When a metric proved unable to resolve the
intended signal, the change is logged inline with a timestamp in the same
document rather than applied silently: the $\tau{=}0.99\to0.9$ knee-metric
re-registration in \S\ref{sec:taskd-m3} happened \emph{before} the
force-rank sweep launched, triggered by a calibration run showing the
reliably-trainable model plateaus at unconstrained cosine
$\approx0.936$ -- a ceiling that would have made $\tau{=}0.99$
uninformative for every cell, discovered by deliberately checking the
easiest case first rather than by post-hoc metric selection after seeing
the sweep's results.

\subsection{The adversarial gauntlet}
\label{sec:appendix-gauntlet}

Before any experiment's code ran on the cluster, design and code passed
through four stages, each performed or reviewed by a fresh context that
did not write the artifact under review -- this project's standing rule
that an implementer does not audit their own work:

\begin{enumerate}
\setlength\itemsep{2pt}
\item \textbf{Task/design attack.} A pass whose only job is to find a
      shortcut that would make a positive result meaningless. This stage
      killed the program's original design outright before any training
      code was written: an early $K$-parallel-entity design (Task
      A$'$) admitted the rank-1 ``vector-in-a-costume'' shortcut
      $Z=uv_0^\top$ (\texttt{ATTACK\_task\_shortcuts.md}), and a
      separate pass found the proposed crossover prediction was not
      interpretable even if observed (\texttt{ATTACK\_baseline\_
      confound.md}). Both forced the pivot to the provable-bound,
      pinned-readout construction used throughout this paper
      (\S\ref{sec:setup}). The same stage, applied to the Task E design,
      is the source of the general-permutation periodicity finding in
      \S\ref{sec:taske}'s footnote.
\item \textbf{Novelty check.} A literature pass confirming the specific
      claim -- not the general area -- had not already been shown; see
      \S\ref{sec:setup-bound} and \S\ref{sec:related} for the resulting
      scoping of what is and is not claimed as new.
\item \textbf{Code audit}, typically 2--3 independent rounds per
      experiment. For Task E, the first round
      (\texttt{AUDIT\_task\_e\_correctness.md},
      \texttt{AUDIT\_task\_e\_validity.md}) caught two FATAL-severity
      bugs that would have silently invalidated the headline result: an
      injectivity check on the key$\to$value mapping carried a
      floating-point tolerance that could not detect a single merged
      binding (defeating the $\rank(Z)\geq K$ guarantee for exactly the
      instances it was supposed to catch), and a hop-depth generator that
      sampled a general random permutation rather than a single
      Hamiltonian $K$-cycle, so held-out hops silently collapsed into
      in-distribution or trivial queries via $h\bmod(\text{cycle
      length})$ (measured up to 100\% collapse at small $K$ before the
      fix). Both fixes were independently re-audited fresh
      (\texttt{AUDIT\_task\_e\_reaudit.md}) rather than trusted from a
      self-check -- a second round that itself caught residual issues
      the first fix's self-verification had missed, the concrete reason
      this project's process treats one audit round as insufficient.
\item \textbf{Runtime smoke gates}, required to pass immediately before
      the first real training step: (i) the gradient-based blank-out
      test (\S\ref{sec:setup-bottleneck}) -- corrupt the raw-input cache
      after the encoder writes $Z$ and confirm the decode path is
      bit-for-bit unchanged, the only evidence treated as sufficient that
      the single-state bottleneck actually holds; (ii) a degenerate-
      spectrum backward-pass check on the \texttt{eigh}-based
      effective-rank computation, since eigendecomposition backward
      passes are numerically unstable at repeated eigenvalues -- a
      regime the orthonormal-key construction (which gives the exact
      target memory a perfectly degenerate spectrum, all singular values
      equal, \S\ref{sec:setup-task}) deliberately induces; (iii) a
      resolution pre-flight confirming the exact closed-form memory
      $Z=\sum_jv_jk_j^\top$ reproduces itself under the harness's own
      scoring code before any learned model is trained against it.
\end{enumerate}

\subsection{The extended-budget rule is a discovered process artifact, not
an assumption}
\label{sec:appendix-budget}

The rule stated and applied in \S\ref{sec:frontier-pattern} --
re-test a cell at $2$--$2.5\times$ its original step budget before
calling it dead -- was not decided in advance; it was extracted after the
same false-negative pattern recurred three independent times across three
different axes of this program (hyperparameter, task, and model-scale
axes respectively; enumerated in \S\ref{sec:frontier-pattern}). We report
the rule together with its counter-example (\S\ref{sec:frontier-pattern}'s
$d{=}32$ genuine plateau, which extended budget does \emph{not} rescue)
specifically so it is not mistaken for a license to keep re-testing
indefinitely until a cell turns positive -- a form of the same optional-
stopping error the pre-registration discipline in
\S\ref{sec:appendix-prereg} is designed to prevent elsewhere in this
program.

\section{Supplementary Figure: An Eval-Time Truncation Staircase}
\label{sec:appendix-supp-figure}

\begin{figure}[h]
\centering
\includegraphics[width=0.95\linewidth]{fig_eval_trunc_staircase.pdf}
\caption{\textbf{[Supplementary / illustrative -- not one of this paper's
pre-registered claims and not behind any table in the main text.]} A
different causal instrument from this paper's primary test
(Figure~\ref{fig:taskd-forcerank}, which forces rank at \emph{train} time):
here a single unconstrained-rank $d{=}64,K{=}32$ operator is trained once,
then truncated to rank $k$ \emph{post hoc at eval time} via SVD, and
re-scored at several hop depths. We include it only because it is a
strikingly clean picture of the same underlying mechanism -- a knee at
$k{=}K$ that sharpens with composition depth $h$ -- at an operating point an
order of magnitude larger than Table~\ref{tab:taskd-m3}'s cells. It is
\emph{not} evidence for this paper's causal claim on its own terms: Task D's
own M2 eval-time-truncation probe was informative there only because M3's
train-time force-rank result was pre-registered as primary
(\S\ref{sec:setup-arch}); post-hoc truncation alone was the uninformative
instrument in the companion paper's CODI setting (\S\ref{sec:intro}). Read
this figure as a mechanism illustration, not as a fourth confirmed cell of
Table~\ref{tab:frontier-d4} -- its $(d,K)=(64,32)$ operating point and step
budget do not match any row of that table.}
\label{fig:eval-trunc}
\end{figure}
